\documentclass[10pt, leqno]{article}

\usepackage{diagbox}
\usepackage{array}
\usepackage{float}

\begin{document}
    \subsubsection*{A1}
    Ein \textbf{Programm} ist eine kodierte Folge von Anweisungen die, nach eventueller Übersetzung, von einem Prozessor ausgeführt werden können.
    \\\\
    Ein \textbf{Prozess} ist ein vom Betriebssystem instanziertes Programm. Ein Prozess verfügt über Eigenschaften, die es dem Betriebssystem ermöglichen effizient die Resourcen der Hardware auf die laufenden Prozesse aufzuteilen.
    \\\\
    Der Begriff \textbf{Thread} bezeichnet zwei sich ähnelnde Konzepte.
    \begin{itemize}
        \item \textbf{Logische Threads} dienen zur parallelisierung innerhalb eines Prozess.
        So können mehrere Abschnitte eines Programms gleichzeitig ausgeführt werden.
        Der Scheduler eines Betriebssystems verteilt die Resourcen der Hardware auf diese Threads.
        Ein Prozess kann aus beliebig vielen Threads bestehen. 
        \item \textbf{physische Threads} sind die tatsächlichen Threads, die ein Prozessorkern parallel ablaufen lassen kann.
        Die meisten modernen Prozessoren unterstützen zwei Threads pro Kern.  
    \end{itemize}
    \subsubsection*{A2}
    Das Amdahlsche Gesetz berechnet den Speedup $\eta_S$ durch:
    $$\eta_S = \frac{T}{t_s + \frac{t_p}{n_p}}$$
    Es ergibt sich die folgende Tabelle
    
    \begin{table}[H]
        \centering
        \begin{tabular}{c | c | c | c | c}
            \diagbox{$t_p = $}{$2x$ Kerne} & $x = 0$ & $x = 1$ & $x = 2$ & $x = 3$\\\hline
            0.25 & undef. & 1.14 & 1.23 & 1.26 \\\hline
            0.5  & under. & 1.33 & 1.60 & 1.71 \\\hline
            0.75 & undef. & 1.60 & 2.29 & 2.67 \\\hline
        \end{tabular}
    \end{table}
    
    Könnte es sein, dass $2^x$ gemeint war in der Aufgabe? Dann sähe die Tabelle wie folgt aus
    
    \begin{table}[H]
        \centering
        \begin{tabular}{c | c | c | c | c}
            \diagbox{$t_p = $}{$2^x$ Kerne} & $x = 0$ & $x = 1$ & $x = 2$ & $x = 3$\\\hline
            0.25 & 1.00 & 1.14 & 1.23 & 1.28 \\\hline
            0.5 & 1.00 & 1.33 & 1.60 & 1.78 \\\hline
            0.75 & 1.00 & 1.60 & 2.29 & 2.91 \\\hline
        \end{tabular}
    \end{table}
    
    \subsubsection*{A3}
    


\end{document}